\chapter{Relación con la Asignatura Programación 1 (INF120)}\label{cap:inf120}
Este capítulo muestra que el proyecto no es un tema ajeno a la asignatura, sino una aplicación integradora de sus contenidos. Cada concepto de Programación 1 ---tipos de datos, estructuras de control, ciclos, funciones, arreglos, matrices, registros, modularidad, manejo de errores y entrada de argumentos--- aparece en el código fuente cumpliendo una función concreta dentro del sistema que aprende. La exposición sigue el orden habitual del programa de la asignatura y cada sección se ilustra con un fragmento extraído automáticamente del repositorio, de modo que el texto y el código no puedan divergir.

\section{Correspondencia entre contenidos y código}
La Tabla~\ref{tab:inf120} resume la correspondencia entre los contenidos de \Sigla{} y los elementos del proyecto. La columna central indica la construcción de Rust utilizada y la columna derecha, el lugar del código donde puede observarse.

\begin{table}[!htbp]
\caption{Correspondencia entre contenidos de Programación 1 y el código del proyecto}\label{tab:inf120}
\small
\begin{tabularx}{\textwidth}{@{}>{\raggedright}p{3.6cm} >{\raggedright}X >{\raggedright\arraybackslash}p{4.3cm}@{}}
\toprule
Contenido de \Sigla & Construcción en Rust & Ubicación en el proyecto \\
\midrule
Variables y tipos de datos & \texttt{f32}, \texttt{usize}, \texttt{char}, \texttt{\&str}; inmutabilidad por defecto y \texttt{let mut} & Todos los módulos \\
Operadores y expresiones & Aritméticos, relacionales, lógicos; \texttt{exp()}, \texttt{max} & \texttt{Activacion}, softmax \\
Estructuras condicionales & \texttt{if}/\texttt{else}, \texttt{match} con patrones & \texttt{main.rs}, \texttt{parse\_vowel} \\
Ciclos & \texttt{for} sobre rangos, \texttt{for} con \texttt{enumerate}, iteradores & Épocas y muestras de cada experimento \\
Funciones y parámetros & Paso por valor y por referencia (\texttt{\&}, \texttt{\&mut}) & \texttt{forward}, \texttt{update\_weights} \\
Arreglos y vectores & \texttt{Vec<T>}, rebanadas \texttt{\&[f32]}, arreglos fijos \texttt{[[\&str; 5]; 5]} & \texttt{Vector<T>}, patrones de vocales \\
Matrices & Almacenamiento por filas; patrones $5\times5$ aplanados & \texttt{Matrix<T>}, \texttt{encode} \\
Registros & \texttt{struct} con métodos en bloques \texttt{impl} & \texttt{Neuron}, \texttt{Layer}, \texttt{Network} \\
Genéricos & \texttt{Vector<T>} con restricciones de traits & \texttt{math/} \\
Modularidad & \texttt{mod}, \texttt{pub}, \texttt{use} & \texttt{math/}, \texttt{neural/}, \texttt{learning/}, \texttt{experiments/} \\
Validación y errores & \texttt{assert\_eq!}, \texttt{Result<T, E>}, \texttt{Option} & Dimensiones, vocal ingresada \\
Pruebas y depuración & Trazas por época, prueba de escritorio, \texttt{\#[test]} & Capítulo~\ref{cap:implementacion}, motor complementario \\
\bottomrule
\end{tabularx}
\fuente{Elaboración propia a partir del código fuente de \url{\Repositorio}.}
\end{table}

\section{Tipos de datos, variables y expresiones}
Toda la aritmética de la red del repositorio oficial se realiza con números de punto flotante de precisión simple (\texttt{f32}, IEEE~754), que representan unas siete cifras decimales significativas. Es una elección adecuada para una demostración: los tres experimentos alcanzan su objetivo, y el error final de la regresión $y=2x$ (\OfLinMseFin) se acerca al límite de esa precisión. Los índices y tamaños usan \texttt{usize}, el entero sin signo del tamaño de una dirección de memoria, y la vocal ingresada se maneja como \texttt{char}. En Rust las variables son inmutables por defecto; los pesos que cambian durante el entrenamiento se declaran con \texttt{let mut}, lo que hace visible en el código qué datos evolucionan.

\section{Funciones, parámetros y referencias}
Las operaciones matemáticas se implementan como funciones con parámetros tipados y valor de retorno explícito. El Código~\ref{cod:dot} muestra el producto punto genérico de \texttt{Vector<T>}: recibe el otro vector por referencia inmutable (\texttt{\&Vector<T>}), comprueba que las longitudes coincidan y acumula los productos con iteradores. Pasar por referencia evita copiar los datos, y el compilador garantiza que la función no los modifique.

\lstinputlisting[caption={Producto punto genérico con validación de longitudes (\texttt{src/math/vector.rs})},label=cod:dot]{generado/codigo/of_dot.rs}

\section{Estructuras de control}
El punto de entrada (Código~\ref{cod:main}) decide qué experimento ejecutar con un \texttt{match} sobre el primer argumento de la línea de comandos. A diferencia de un \texttt{switch} tradicional, el compilador exige que \texttt{match} cubra todos los casos posibles; el comodín \texttt{Some(value)} trata cualquier otro texto como una vocal y \texttt{None}, la ausencia de argumentos.

\lstinputlisting[caption={Selección del experimento con \texttt{match} (\texttt{src/main.rs})},label=cod:main]{generado/codigo/of_main.rs}

El Código~\ref{cod:argmax} busca el índice del valor máximo, operación que convierte las cinco probabilidades de la red de vocales en una clase. El algoritmo es el clásico de Programación 1, expresado con iteradores: recorre los pares (índice, valor) y conserva el mayor.

\lstinputlisting[caption={Búsqueda del índice máximo (\texttt{src/experiments/vowels.rs})},label=cod:argmax]{generado/codigo/of_argmax.rs}

\section{Ciclos}
El entrenamiento es un conjunto de ciclos anidados: un ciclo externo sobre las épocas y un ciclo interno sobre las muestras. En el experimento $y=2x$ (Código~\ref{cod:of-lin}, Capítulo~\ref{cap:implementacion}) el ciclo externo \texttt{for epoch in 1..=max\_epochs} se interrumpe con \texttt{break} cuando el error alcanza la meta, y el interno acumula los gradientes. El mismo patrón, con \texttt{enumerate} para conocer el índice de cada muestra, aparece en los experimentos booleano y de vocales.

\section{Arreglos, vectores y matrices}
La estructura \texttt{Vector<T>} encapsula un arreglo dinámico \texttt{Vec<T>}, y \texttt{Matrix<T>} almacena sus $m\times n$ elementos en un único arreglo contiguo en orden por filas: el elemento de la fila $i$ y columna $j$ ocupa la posición $i\cdot n+j$. El Código~\ref{cod:mulvec} muestra la multiplicación de matrices con tres ciclos anidados, la implementación directa de $c_{ij}=\sum_k a_{ik}b_{kj}$.

\lstinputlisting[caption={Multiplicación de matrices sobre almacenamiento por filas (\texttt{src/math/matrix.rs})},label=cod:mulvec]{generado/codigo/of_multiply.rs}

Los arreglos de tamaño fijo también aparecen en el proyecto. Las cinco vocales se dibujan en un arreglo \texttt{[[\&str; 5]; 5]}, donde \texttt{\#} representa un píxel encendido (Código~\ref{cod:patrones}). La función \texttt{encode} (Código~\ref{cod:encode}) recorre la matriz fila por fila y la aplana en un vector de 25 valores 0 o 1: el mismo recorrido por filas que se enseña para matrices en la asignatura.

\lstinputlisting[caption={Patrones binarios $5\times5$ de las vocales (\texttt{src/experiments/vowels.rs})},label=cod:patrones]{generado/codigo/of_patrones.rs}
\lstinputlisting[caption={Aplanado de un patrón en 25 entradas (\texttt{src/experiments/vowels.rs})},label=cod:encode]{generado/codigo/of_encode.rs}

\section{Registros y procedimientos}
Los registros de Programación 1 se corresponden con las estructuras (\texttt{struct}) de Rust, y los procedimientos que operan sobre ellos, con los métodos de los bloques \texttt{impl}. La neurona (Código~\ref{cod:neuron-struct}) es un registro con dos campos privados, sus pesos y su bias; el exterior solo puede usarla mediante sus métodos públicos \texttt{new} y \texttt{forward}, lo que aplica el principio de encapsulamiento.

\lstinputlisting[caption={Registro de una neurona (\texttt{src/neural/neuron.rs})},label=cod:neuron-struct]{generado/codigo/of_neuron_struct.rs}

\section{Validación de datos y manejo de errores}
Un error frecuente en los primeros cursos es combinar arreglos de tamaños distintos. El repositorio oficial lo previene con \texttt{assert\_eq!}: si las dimensiones no coinciden, el programa se detiene con un mensaje que describe la causa, en lugar de producir un resultado incorrecto. Para los datos ingresados por el usuario utiliza \texttt{Result} (Código~\ref{cod:parsear}): la vocal se normaliza a mayúscula y, si no es A, E, I, O o U, se devuelve un error que el programa informa sin interrumpirse de forma abrupta.

\lstinputlisting[caption={Validación de la vocal ingresada (\texttt{src/experiments/vowels.rs})},label=cod:parsear]{generado/codigo/of_parse.rs}

\section{Pruebas, depuración y prueba de escritorio}
La prueba de escritorio, técnica clásica de Programación 1 para seguir manualmente un algoritmo, se aplica en dos niveles. El repositorio oficial imprime en cada época los gradientes y los parámetros, de modo que cada actualización puede comprobarse con calculadora (Tabla~\ref{tab:of-bool}); el motor complementario registra además las primeras iteraciones del entrenamiento (Sección~\ref{sec:escritorio}) e incluye 52 pruebas automatizadas con \texttt{\#[test]} y \texttt{assert!}. Así, verificar un programa se muestra como parte de su construcción.

\begin{nota}[Propuesta de uso en aula]
En \Sigla, los estudiantes pueden reproducir la prueba de escritorio, modificar la tasa de aprendizaje o las épocas, dibujar una nueva letra en los patrones o añadir una función a \texttt{Activacion} con su prueba \texttt{\#[test]}.
\end{nota}
